% 卷首.tex
% Auto-generated from 感悟论语.odt by tools/odt_to_tex.py

\chapter{作者简介}

本书作者田冀焕于 2000 年考入清华大学电机系，在此完成了本科学业、博士与博士后的研究工作。爱好传统文化，坚持践行论语的基本观点与思想。目前，作者独立从事高性能数值算法的理论研究与软件开发工作，是开源软件项目 HierBEM 的项目负责人。

\chapter{前言}

我在清华本科期间曾加入旨在弘扬中国传统文化的钟铎社。那时的下午课后，我常与同学在十食堂西侧的小树林内一起朗读朱熹的《四书章句集注》，因而对论语略有涉猎。不过，我当时的阅读与学习并不系统，理解浮于表面、难有心得，更不用说在生活中贯彻践行了。直至2015年，我有幸聆听了傅佩荣老师的《论语300讲》，才算是对论语有了全面的了解和一定的体会。通过持续的学习与践行，我不仅在待人接物方面从中获益匪浅，更是促成了对个人人生方向的探索与逐渐明晰，以及对此方向坚韧不拔的坚守与追寻。

目前市面上常见的与论语相关的书籍大多是针对原文的注释与解读，而展现论语与现代日常生活各方面的诸多关联，以及讲述如何具体应用与践行论语基本原则的书籍较少。本书分享我在学习、科研、职场工作中与论语深度互动、学而时习的体会感悟、心路历程，希望“以文会友，以友辅仁”，向广大读者提供“与孔子对话”的阅读体验，给人生带来启发与积极向上的推动作用。目标读者群主要是对传统文化有兴趣、并希望通过传统文化改善与提升自己的认知与境界的科研工作者、学者、学生、普通读者以及职场人。

目前书稿篇幅约九万余字，包括修养、科研、职场、习惯、家庭、社会、论语解读七部分。每篇文章独立成章，既可连续阅读，也能随手翻阅。此外，在附录中整理了论语章句分类以及孔子与弟子的对话，方便读者系统查阅。
